\part{Geometric reconstruction}
\chapter{What the comparison maps preserve}
Extensions supply Hall multiplication through their flag correspondences.
The cyclic pairing records duality on the underlying category.
A coproduct combines representations, and a half-braiding specifies
central exchange with every object.
These operations have distinct carriers.
A reconstruction must construct the maps between them before comparing
traces or characters.

For a specified primitive Lie map, Theorem~\ref{cyiv:primitive-comparison}
constructs its enveloping and Chevalley extensions explicitly.
Surjectivity and equality of finite parity-refined PBW components then
prove an algebraic isomorphism.
For a specified strong monoidal equivalence,
Theorem~\ref{cyarch:centre-transport} constructs the transported
half-braiding and proves a braided equivalence of categorical centres.
Neither statement constructs a geometric functor from a scalar character.

\chapter{K3 actions and their common carriers}
The Mukai lattice has rank twenty-four and symmetric pairing of signature
$(4,20)$.
The Nakajima operators act on the Hilbert-scheme cohomology tower.
Lefschetz operators act on each selected hyperk\"ahler moduli space.
The local ADE Yangian actions use their equivariant surface coefficients.
These state spaces and coefficient fields are part of the respective
constructions.

On smooth projective moduli spaces, Proposition~\ref{cyiv:k3-kernels}
identifies rational cohomology kernels with graded linear operators and
convolution with composition.
A chosen collection of geometric operators therefore gives the faithful
action algebra of Proposition~\ref{cyiv:k3-action}.
Its defining ideal consists of the actual relations of those operators.
The construction neither inserts a spectral coproduct nor identifies
local equivariant actions with global ones.
Those comparisons require intertwining maps on their state spaces.

The coend of Proposition~\ref{cyiv:k3-coend} instead constructs matrix
coefficients for a specified tensor category and fibre functor.
The selected kernels become comodule morphisms.
This separates the algebra generated by geometric operations from tensor
symmetries commuting with them.
A Yangian realization must construct its current generators, relations,
and tensor action on these geometric carriers.

\chapter{The compact Hall and automorphic comparison}
For $X=S\times E$, the $E$-equivariant motivic Hall construction retains
extensions of sheaves and the reduced integration map on its regular
subalgebra.
A chosen charge homomorphism gives the polarized character of
Proposition~\ref{cyiv:polarized-character}.
The Borcherds algebra and its root Fock character are independently defined.

A primitive root comparison must preserve brackets, parity, and charge.
Its graph quotient is computed in Proposition~\ref{cyiv:graph-quotient}.
The quotient is the geometric Borel for every specified comparison;
it is the Borcherds Borel as well only when that comparison is an
isomorphism.
The quotient construction does not provide the comparison map.

A geometric realization also needs the critical coefficient complexes
and extension-compatible square roots of their determinant lines.
The cross-line $\det R\operatorname{Hom}(A,B)$ enters the orientation
formula for an extension of $B$ by $A$.
Compatibility on triple flags makes this data usable in Hall convolution.
The brane functor, paired quantum coproduct, and trace-line comparison
remain the maps specified in Problem~\ref{cyiv:compact-comparison}.

\chapter{Classical fields, centres, and amplitudes}
The cyclic $A_\infty$ action satisfies the classical master equation by
Proposition~\ref{cyiv:cyclic-master}.
The K3 field reduction of Proposition~\ref{cyiv:k3-classical-reduction}
retains both the degree-zero and degree-two Dolbeault modes.
The finite one-loop Lie weights vanish by their positive-height trace
calculation.
Analytic quantization and compatible observable completions are separate
steps in a quantum field comparison.

For an ordinary enveloping algebra, Theorem~\ref{cyiv:hh-ce} computes the
Hochschild centre complex using Chevalley cochains with adjoint
coefficients.
Its explicit antisymmetrization map identifies the two ordinary
cochain models.
A physical bulk comparison must construct its map into this centre or
its chosen geometric refinement, preserving the actual local operations.
A state and its pairing must then be compared separately.
Thus the algebraic calculations specify concrete parts of reconstruction
and the geometric maps still required to recover amplitudes.
